\chapter*{Resumen}

\section*{\tituloPortadaVal}

PITEA es una herramienta en línea de comandos que permite la ocultación de archivos de texto mediante el cifrado de estos, su posterior ocultación o transformación en imágenes y, finalmente, en audio. También permite realizar el proceso inverso. La aplicación combina técnicas de criptografía y esteganografía para ocultar la información que pueda ser transmitida en la red pública, aprovechando la naturaleza digital de las imágenes y el sonido. 
A través de  nuestra herramienta, los usuarios podrán garantizar la confidencialidad de sus datos, eligiendo entre distintos métodos de ocultación. El objetivo es desarrollar una herramienta versátil y accesible para cualquier persona con un mínimo conocimiento del uso de la línea de comandos, permitiendo proteger información sensible “a plena luz del día”. Este enfoque amplía las aplicaciones de la criptografía y la esteganografía, ofreciendo nuevas formas de ocultar y proteger información digital.



\section*{Palabras clave}
   
\noindent Criptografía, esteganografía, línea de comandos, ocultación, cifrado, red pública, confidencialidad, información sensible, protección de datos.

   


